\chapter{Declarations}
\label{ch:declarations}

\section{Use of Generative AI}
I acknowledge the use of Claude and Claude Code (Anthropic, https://claude.ai/) and the underlying models (Opus 4.8 and Sonnet 4.6) while undertaking this work. I wrote the reimplementation of the Contrastive NCM detector that this thesis evaluates and all of the attack implementations. AI was used to write supporting code that is not part of the thesis's contribution: the multi-layer-perceptron base-learner module, the CICIDS2017 data-loading and preprocessing utilities, and the figure- and table-generation helpers. Under my own design and supervision, it also assisted with implementing the streaming harness and I used it for boilerplate in the experiment notebooks and for debugging and refactoring throughout. For the report itself, the content, findings, and contributions (the threat model, the attack taxonomy, the experimental design, and the analysis and conclusions) are my own. I primarily used AI for proofreading, refining language, and structured feedback on drafts. All AI-generated code was read, tested, and integrated by me, and I am responsible for its correctness. No AI output has been presented as my own original contribution or relied upon as a factual source without independent verification.

\section{Ethical Considerations}
This work takes an adversarial perspective on a deployed adaptive intrusion-detection architecture, developing concrete attacks that force or suppress its adaptation and poison or backdoor it through its retraining set~\cite{Kuppa_2022}. Therefore, it falls under \emph{offensive security} or \emph{red teaming}. Documenting such attacks carries the inherent risk of providing a blueprint for adversaries, but ``security through obscurity'' is a fragile defence, and it is preferable for these weaknesses to be identified and documented here than discovered by real-world attackers. To that end, every attack is reported alongside the mechanism that enables it (\cref{ch:evaluation,ch:conclusion}).

All experiments are conducted entirely within a sandboxed simulation, on local or authorised university hardware, using only the publicly available CICIDS2017 benchmark. No live systems, third-party services, or commercial APIs are probed or targeted, and no personally identifiable information is involved. The project is conducted in accordance with the \emph{Computer Misuse Act 1990 (UK)}.

\section{Sustainability}
The computational footprint of this work is modest and was actively managed. Experiments were run on a single NVIDIA RTX 3090 GPU (a rented \texttt{vast.ai} instance) and on an Imperial-provided virtual machine with a Tesla T4 GPU. The dominant cost is the repeated retraining of the detector's encoder and the base learner at each concept-discovery event. This is bounded by the relatively small models involved and by fixed, short training schedules rather than open-ended optimisation. The streaming harness keeps data GPU-resident to avoid repeated host--device transfers, and design choices were fixed from a small number of targeted runs rather than exhaustive hyperparameter searches. Relative to contemporary deep-learning workloads the datasets and models used here are small, so the overall energy use and carbon footprint of the project are correspondingly low.

\section{Availability of Data and Materials}
This work relies exclusively on publicly available data. The evaluation uses the CICIDS2017 intrusion-detection benchmark~\cite{Sharafaldin_2018} in its Engelen-corrected release~\cite{Engelen_2021}, both of which are openly available to researchers. No proprietary or personally identifiable information was collected, processed, or generated, so the risk of privacy harm is negligible. Because Kuppa and Le-Khac~\cite{Kuppa_2022} publish no reference implementation, the Contrastive NCM detector, the coupled adaptive system, and all attacks were reimplemented from scratch for this project. The full source code, together with the experiment notebooks, is available at https://github.com/dillan-scott/final-year-project-dbs21.